\section{Introduction}\label{sec:introduction}

Cutting a quantum field theory into spatial regions replaces a single system by theories with artificial boundaries. Once boundary conditions are chosen, each region can be quantized independently, giving regional Hilbert spaces and Hamiltonians. A separate question then arises: given these already quantized regional theories, how can one couple them so that the quantum theory on the joined region is recovered?

In this paper we study this question for free scalar fields using a quadratic interface interaction. We begin with two regional theories with Neumann conditions at the artificial cut, quantize them separately, and then couple their boundary field traces through a quadratic interface interaction proportional to the square of their mismatch. The corresponding continuum finite-coupling problem provides a simple analytic interpolation: at finite coupling the interface permits a nonzero trace mismatch, while the strong-coupling limit suppresses that mismatch and recovers the smooth joined scalar. We use this continuum problem as the spectral target for the regulated quantum construction below. This prescription is therefore a concrete realization of gluing for this chosen class of regional theories.

The main complication appears after regional quantization is regulated. A mode truncation removes high-frequency oscillators, but these omitted modes still contribute to the response of the boundary trace. Consequently, using the continuum interface coupling unchanged at finite cutoff does not reproduce the continuum interface problem exactly, even at low energy. Likewise, imposing trace equality directly as a hard constraint on the retained modes ignores the boundary response carried by the discarded sector. Thus the ultraviolet truncation affects the infrared gluing condition through the boundary response.

We compensate for this effect by matching the boundary resolvent of the truncated theory to that of the continuum problem at a chosen spectral parameter. Equivalently, the contribution of the omitted modes is absorbed into a cutoff-dependent interface coupling. For the interval example this running coupling can be obtained analytically. At the perfectly joined endpoint it diverges with the cutoff, as expected for a penalty enforcing trace continuity, but its coefficient is fixed by the omitted static response rather than chosen ad hoc. In a fixed low-energy spectral window, the leading response-matched coupling improves the convergence of the joined spectrum from the \(O(N^{-1})\) behavior of a direct hard projection to \(O(N^{-2})\).

Recovering the spectrum is not sufficient to reconstruct the quantum theory. The ground-state covariance also depends on the eigenvectors of the regulated Hamiltonian and on how observables are represented in the truncated regional spaces. We therefore use the same response-matched Hamiltonian, without any additional fitting, to test selected equal-time vacuum two-point functions. For two Neumann half intervals we consider both a spacelike correlation across the interface and one-sided interface-to-bulk correlations defined by shrinking smooth packets. In the tested cutoff sequence, the regulated glued covariances converge to independent joined-theory references, including agreement of the two one-sided interface limits after the known finite-width asymmetry of the packets is removed.

The interval is the main example because both the continuum interface problem and the regulated regional theory are explicitly controllable. Two supplementary tests probe aspects that are independent of the one-dimensional trigonometric spectrum. A four-region square checks that the same local interaction can be assembled edge by edge and partially or completely glued, while a global \(\mathrm{AdS}_2\) example tests the same response-matching mechanism in a curved background.

The scope of these results is deliberately limited. We establish a controlled interface interaction construction for the displayed free Gaussian models, together with fixed-window spectral convergence and convergence of selected cutoff-dependent smeared vacuum correlators. We do not prove convergence for arbitrary states or observables, a unitary equivalence between regional and joined continuum Fock representations, a regulator-uniform ultraviolet theorem, or an interacting-QFT gluing construction.

Section 2 develops the regional scalar theory, its covariant phase space and canonical quantization, the interface interaction, and the boundary-response matching rule. Section 3 applies the construction to two intervals, from the exact finite-coupling spectrum through the regulated Hamiltonian and running coupling to the spectral and vacuum-correlation tests. Section 4 summarizes the conclusions. Appendix A contains the four-region compositional check, and Appendix B gives the global \(\mathrm{AdS}_2\) check.
